% FILE: CSE-19_CT-2.tex
\documentclass[11pt, a4paper]{article}
\usepackage[a4paper, top=2.5cm, bottom=2.5cm, left=2cm, right=2cm]{geometry}
\usepackage{fontspec}
\usepackage[english]{babel}
\usepackage{amsmath}
\usepackage{amssymb}
\usepackage{enumitem}

\babelfont{rm}{Noto Sans}

\begin{document}

%%%%%%%%%%%%%%%%%%%%%%%%%%%%%%%%%%%%%%%%%%%%%%%%%%  
% QUESTION_START  
%%%%%%%%%%%%%%%%%%%%%%%%%%%%%%%%%%%%%%%%%%%%%%%%%%
% id=cse19-ct2-seca-q1  
% discipline=CSE  
% year=19  
% exam_type=CT  
% section=A  
% ct_number=2  
% question_number=1  
% topic=Differentiation  
% subtopic=Leibnitz Rule  
% difficulty=Hard  
% length=Long  
% question_type=New  
% frequency=1  
% appearances=  
% OCR_STATUS=VERIFIED

\section*{CT-2 (Sec A) - Q1}

Question:
State Leibnitz's theorem. Find the value of the $n$-th derivative of $y = e^{m\sin^{-1}x}$ at $x = 0$.

Solution:

\begin{description}
    \item[Step 1: State Leibnitz's Theorem] \ \\
    If $u$ and $v$ are two functions of $x$ possessing derivatives of the $n$-th order, then the $n$-th order derivative of their product $uv$ is given by:
    \[ (uv)_n = u_n v + \binom{n}{1} u_{n-1} v_1 + \binom{n}{2} u_{n-2} v_2 + \dots + \binom{n}{r} u_{n-r} v_r + \dots + u v_n \]
    where subscripts represent the order of differentiation.

    \item[Step 2: Find successive derivatives of $y = e^{m\sin^{-1}x}$] \ \\
    Given:
    \[ y = e^{m\sin^{-1}x} \implies y(0) = e^0 = 1 \]
    Differentiating once with respect to $x$:
    \[ y_1 = e^{m\sin^{-1}x} \cdot \frac{m}{\sqrt{1-x^2}} = \frac{my}{\sqrt{1-x^2}} \implies y_1(0) = m \]
    Squaring both sides to remove the radical:
    \[ y_1^2 (1-x^2) = m^2 y^2 \]
    Differentiating again with respect to $x$:
    \[ 2y_1 y_2 (1-x^2) + y_1^2 (-2x) = m^2 (2y y_1) \]
    Dividing through by $2y_1$:
    \[ y_2(1-x^2) - xy_1 = m^2 y \]
    At $x = 0$:
    \[ y_2(0) - 0 = m^2 y(0) \implies y_2(0) = m^2 \]

    \item[Step 3: Apply Leibnitz's Theorem to get a recurrence relation] \ \\
    Differentiating the equation $y_2(1-x^2) - xy_1 - m^2 y = 0$ $n$ times using Leibnitz's theorem:
    \[ \left[ y_{n+2}(1-x^2) + n y_{n+1}(-2x) + \frac{n(n-1)}{2} y_n (-2) \right] - \left[ y_{n+1}x + n y_n \right] - m^2 y_n = 0 \]
    \[ (1-x^2)y_{n+2} - (2n+1)xy_{n+1} - (n^2 + m^2) y_n = 0 \]
    At $x = 0$, this simplifies to:
    \[ y_{n+2}(0) = (n^2 + m^2)y_n(0) \]

    \item[Step 4: Find the general value of $y_n(0)$] \ \\
    We use the recurrence relation $y_{n+2}(0) = (n^2 + m^2)y_n(0)$ to find the pattern:
    \begin{itemize}
        \item \textbf{Case 1: If $n$ is even} \\
        Let $n = 2k$ for $k \ge 1$:
        \[ y_2(0) = m^2 \]
        \[ y_4(0) = (2^2 + m^2)y_2(0) = m^2(2^2 + m^2) \]
        \[ y_6(0) = (4^2 + m^2)y_4(0) = m^2(2^2 + m^2)(4^2 + m^2) \]
        In general, for even $n$:
        \[ y_n(0) = m^2(2^2 + m^2)(4^2 + m^2) \dots ((n-2)^2 + m^2) \]

        \item \textbf{Case 2: If $n$ is odd} \\
        Let $n = 2k+1$ for $k \ge 0$:
        \[ y_1(0) = m \]
        \[ y_3(0) = (1^2 + m^2)y_1(0) = m(1^2 + m^2) \]
        \[ y_5(0) = (3^2 + m^2)y_3(0) = m(1^2 + m^2)(3^2 + m^2) \]
        In general, for odd $n$:
        \[ y_n(0) = m(1^2 + m^2)(3^2 + m^2) \dots ((n-2)^2 + m^2) \]
    \end{itemize}
\end{description}

Final Answer:
\[ 
\boxed{y_n(0) = \begin{cases} 
m^2(2^2 + m^2)(4^2 + m^2) \dots ((n-2)^2 + m^2), & \text{when } n \text{ is even} \\ 
m(1^2 + m^2)(3^2 + m^2) \dots ((n-2)^2 + m^2), & \text{when } n \text{ is odd} 
\end{cases}} 
\]

%%%%%%%%%%%%%%%%%%%%%%%%%%%%%%%%%%%%%%%%%%%%%%%%%%  
% QUESTION_END  
%%%%%%%%%%%%%%%%%%%%%%%%%%%%%%%%%%%%%%%%%%%%%%%%%%

\end{document}